\subsection{Local Computation of the Degree}

In this section we proof the following theorem.

\begin{theorem}
    \label{alg_top:thm:local_degree}
    Let $f \colon S^n \to S^n$ be smooth around a point $x \in S^n$, $y = f(x)$ and $T_xf \colon T_xS^n \to T_yS^n$ an isomorphism. Then for the local degree of $f$ at $x$ we have 
    \[
    \mathrm{deg}_x(f) = +1
    \]
    if $T_xf$ is orientation preserving and 
    \[
    \mathrm{deg}_x(f) = -1
    \]
    if $T_xf$ is orientation reversing.
\end{theorem}

\begin{proof}
    First assume $T_xf$ is orientation preserving. We have to show that for an open neighbourhood $U \subseteq S^n$ of $x$ the map 
    \[
        \begin{tikzcd}[column sep = 2.2em]
            {H_n(S^n)} & {H_n(S^n,S^n-y_1)} & {H_n(U,U-x)} & {H_n(S^n,S^n-y_1)} & {H_n(S^n)}
            \arrow["\cong", from=1-1, to=1-2]
            \arrow["\cong"', from=1-3, to=1-2]
            \arrow["\cong","H_n(f|_U)"', from=1-3, to=1-4]
            \arrow["\cong"', from=1-5, to=1-4]
        \end{tikzcd}
    \]
    is the identity (where all other maps are induced by inlcusions). This follows immediately from the following claim.

    \begin{claim}
        For orientation preserving embeddings $g_i \colon U \to S^n$ with $g_i(x) = y_i$, $i = 1,2$ the map 
        \begin{equation}
            \label{alg_top:eq:local_degree_claim}
            \tag{$\star$}
            \begin{tikzcd}[column sep = 2em]
                {H_n(S^n)} & {H_n(S^n,S^n-y_1)} & {H_n(U,U-x)} & {H_n(S^n,S^n-y_1)} & {H_n(S^n)}
                \arrow["\cong", from=1-1, to=1-2]
                \arrow["\cong"',"H_n(g_1)", from=1-3, to=1-2]
                \arrow["\cong","H_n(g_2)"', from=1-3, to=1-4]
                \arrow["\cong"', from=1-5, to=1-4]
            \end{tikzcd}
        \end{equation}
        is the identity
    \end{claim}
    Firstly we may assume, that $y_1 = y_2$, otherwise replace $g_1$ with $R \circ g_1$, where $R \colon S^n \to S^n$ is a rotation with $R(y_1) = y_2$. Then the map (\ref{alg_top:eq:local_degree_claim}) is unchanged as seen by the commutativity of the following diagramm
    \[\begin{tikzcd}[every label/.append style={font=\small}, sep =1.5em]
	{H_n(S^n)} & {H_n(S^n,S^n-y_1)} \\
	&& {H_n(U,U-x)} \\
	{H_n(S^n)} & {H_n(S^n,S^n-y_2)}
	\arrow["\cong", from=1-1, to=1-2]
	\arrow["{H_n(R) = id}"', from=1-1, to=3-1]
	\arrow["{H_n(R)}"', from=1-2, to=3-2]
	\arrow["{H_n(g_1)}"', from=2-3, to=1-2]
	\arrow["{H_n(R \circ g_1)}", from=2-3, to=3-2]
	\arrow["\cong", from=3-1, to=3-2]
\end{tikzcd}\]
Note that $R \simeq id_{S^n}$ since $\mathrm{SO}(n)$ is path-connected. Moreover we may assume $U \cong \R^n$, otherwise shrink $U$ to some $U^\prime \subseteq U$ and excision yields $H_n(U^\prime,U^\prime-x) \xrightarrow{\cong}H_n(U,U-x)$. So if we consider $g_1$ and $g_2$ in suitable charts we get orientation preserving embeddings 
\[
\tilde{g}_1,\tilde{g}_2 \colon \R^n \to \R^n
\]
with $\tilde{g}_1(0) = 0 = \tilde{g}_2(0)$. So by Lemma \ref{alg_top:lem:isotopy_lemma} $\tilde{g}_1$ and $\tilde{g}_2$ are isotopic $\mathrm{rel}\, \{0\}$. In particular we get a homotopy
\[
H \colon U \times I \to S^n
\]
from $g_1$ to $g_2$ with $H((U-x) \times I) \subseteq S^n-y_1$, i.e. a homotopy of \textit{pairs}
\[
g_1,g_2 \colon (U,U-x) \to (S^n,S^n-y_1)
\]
and thus 
\[
H_n(g_1) = H_n(g_2) \colon H_n(U,U-x) \to H_n(S^n,S^n-y_1)
\]
which proofs the claim. It remains to consider the case where $T_xf$ is orientation reversing. For this consider the following commutative diagramm
\[\begin{tikzcd}[column sep = 1.7em]
	{H_n(S^n)} & {H_n(S^n,S^n-x)} & {H_n(U,U-x)} && {H_n(S^n,S^n-y)} & {H_n(S^n)} \\
	&&&& {H_n(S^n,S^n-r(y))} & {H_n(S^n)}
	\arrow["\cong", from=1-1, to=1-2]
	\arrow["\cong"', from=1-3, to=1-2]
	\arrow["{H_n(f|_U)}"', "\cong", from=1-3, to=1-5]
	\arrow["{H_n(r \circ f|_U)}"',"\cong", bend right = 15, from=1-3, to=2-5]
	\arrow["{H_n(r)}","\cong"', from=1-5, to=2-5]
	\arrow["\cong"', from=1-6, to=1-5]
	\arrow["{H_n(r)}","\cong"', from=1-6, to=2-6]
	\arrow["\cong", from=2-6, to=2-5]
\end{tikzcd}\]
where $r \colon S^n \to S^n$ is the reflection along the equator. Then $T_x(r \circ f)$ is orientation preserving and the statement follows from what has been shown already and the fact that $\mathrm{deg}(r) = -1$.
\end{proof}